% Home Network Diagram
% Author: Michael Bugert
% License: CC BY 4.0
\documentclass[border=3cm]{standalone}


\usepackage[utf8]{inputenc}
\usepackage[sfdefault]{noto}
\usepackage[T1]{fontenc}
\usepackage{enumitem}
\usepackage[fixed]{fontawesome5}
\usepackage{etoolbox}
\usepackage{graphicx}
\usepackage{tcolorbox}
\usepackage{tikz}
\usepackage[hidelinks]{hyperref}
\usepackage{geometry}



\tcbuselibrary{skins,raster}
\usetikzlibrary{positioning,matrix,fit,decorations.pathreplacing}

\graphicspath{{img/}}

% colors for subnetworks
\definecolor{subnet-none}{HTML}{000000}
\definecolor{subnet-home}{HTML}{E57400}
\definecolor{subnet-guest}{HTML}{14C300}
\definecolor{subnet-vpn}{HTML}{AD0022}
\definecolor{docker-bright}{HTML}{0A97C4}
\definecolor{docker-dark}{HTML}{088AB4}

\tcbset{
	skin=enhanced,
	boxrule=0.1em,
	left=0.25em,
	right=0.25em,
	% box with hostname and device specs
	host description box/.style={
		width=19em,
		sidebyside,
		righthand ratio=0.4,
		fonttitle=\bfseries\large,
		fontupper=\small,
		fontlower=\small
	},
	% tcbraster for network interfaces
	interfaces raster/.style={
		raster columns=1,
		raster width=6.75em,
		% fix raster height to accomodate for fontawesome icon height differences
		raster height=4.25em,
		raster row skip=0.4em,
		raster every box/.style={
			fontupper=\ttfamily,
			left=0em,
			right=0em,
			top=.15em,
			bottom=.15em,
			valign=center
		}
	},
	% tcbraster for host services
	services raster/.style={
		raster columns=1,
		boxrule=0.1em,
		raster width=14em,
		raster row skip=0.4em,
		raster every box/.style={
			left=0em,
			right=0em,
			top=.15em,
			bottom=.15em,
			fontupper=\small,
			boxrule=0.1em
		}
	},
	% box around host services:
	% Never drawn, only used as a container. Width is set to the same width as "services". If width is not specified, the box would extend much further towards the right because the default width is \linewidth (uncomment 'blankest' to see it).
	services box/.style={
		blankest,
		width=14em
	},
	% box around dockerized host services
	dockerized box/.style={
		%title={Dockerized\hfill\includegraphics[height=12pt]{services/docker}},
		title={Dockerized},
		fonttitle=\bfseries,
		colframe=docker-dark,
		colback=docker-bright,
		left=0.25em,
		right=0.25em
	},
	% box around devices with less details
	devices box/.style={
		width=20em,
		fonttitle=\bfseries\large,
		fontupper=\small,
		attach boxed title to top left,
		boxed title style={frame hidden, opacityback=0.0, left=0em},
		colback=white,
		coltitle=black!75!white
	},
	% tcbraster for devices with less details
	devices raster/.style={
		raster columns=1,
		boxrule=0.1em,
		raster width=20em,
		raster row skip=0.4em,
		raster every box/.style={
			left=0em,
			right=0em,
			top=.15em,
			bottom=.15em,
			fontupper=\small,
			boxrule=0.1em,
			height=3em,
			valign=center
		}
	}
}

\tikzset{
	single icon/.style={
		font={\fontsize{34}{38}\selectfont}
	},
	host description/.style={
	},
	% box drawn around all bits defining a host using tikz fit
	host box/.style={
		inner sep=0.5em,
		draw=black!75!white,
		fill=white,
		rounded corners,
		line width=0.1em
	},
	% box drawn around all bits defining a region using tikz fit
	region box/.style={
		inner sep=2em,
		rounded corners=2em,
		line width=0.1em,
		fill=#1!5!white
	},
	region descriptor/.style={
		outer sep=0,
		execute at begin node={\fontsize{7em}{7em}\selectfont\scshape},
		color=#1!15!white
	},
	% tag for marking devices or services as DNS, DHCP, etc. with a customized color
	tag/.style={
		draw=none,
		fill=#1,
		rectangle,
		rounded corners=2pt,
		outer sep=3pt,
		font=\bfseries\scriptsize,
		text=white
	},
	% connection types
	generic/.style={
		solid,
		line width=0.2em,
		rounded corners
	},
	ethernet/.style={
		generic
	},
	wifi/.style={
		generic,
		dashed,
	},
	vpn/.style={
		generic,
		decoration={
			waves,
			radius=0.8em,
			segment length=.55em,
			angle=30
		}
	}
}

% commands to show device and service icons at fixed size, etc.
\newcommand{\interfaceImage}[1]{%
	\raisebox{-.1\height}{\includegraphics[height=12pt]{#1}}%
}
\newcommand{\deviceImage}[1]{%
	\includegraphics[width=5em, height=3em, keepaspectratio]{devices/#1}%
}
\newcommand{\deviceImageInline}[1]{%
	\raisebox{-.25\height}{%
		\parbox[t]{2em}{%
			\centering%
			\includegraphics[width=2em, height=2em, keepaspectratio]{devices/#1}%
		}%
	}%
}
\newcommand{\ip}[2]{%
	\textbf{\textcolor{#2}{#1}}
}

% clickable URLs shown without protocol
\newcommand{\https}[1]{%
	\href{https://#1}{\nolinkurl{#1}}%
}
\newcommand{\http}[1]{%
	\href{http://#1}{\nolinkurl{#1}}%
}

% usage: \service{image}{textual description}{url}
\newcommand{\service}[3]{%
	\ifblank{#1}{%
		\hspace{2em}%
	}{%
		\raisebox{-.3\height}{%
			\includegraphics[height=2em]{services/#1}%
		}\hspace{.5em}%
	}%
	\parbox[m]{10em}{%
		#2%
		\ifblank{#3}{}{%
			\newline%
			\tiny\textcolor{gray}{#3}%
		}%
	}
}

\newcommand{\planned}{\hfill\scriptsize\textcolor{black!60}{\textit{(planned)}}}
\newcommand{\temporary}{\hfill\scriptsize\textcolor{black!60}{\textit{(temporary)}}}

\begin{document}


% remember picture is needed to retain positions of tcolorboxes (otherwise all edges will be misplaced)


\begin{tikzpicture}[remember picture]
	\pgfdeclarelayer{regions}
	\pgfdeclarelayer{hostboxes}
	\pgfsetlayers{regions,hostboxes,main}
	

	% --------------- start PIO ---------------
	% host description
	\node[host description] (router-description) {%
		\begin{tcolorbox}[
			adjusted title=Ensilage PLC,
			host description box
		]			
		IP: \hspace*{\fill} \ip{X.X.X.X}{subnet-home}
		
			\tcblower
			PFC300
		\end{tcolorbox}
	};
	\node[
		below=0em of router-description.south west,
		anchor=north west
	] (laptop-services) {
		\begin{tcboxedraster}[services raster]{services box}
			\tcbox{\service{}{Application}{}}
		\end{tcboxedraster}
	};
	
	% services
	\node[
		below=0em of router-description.south west,
		anchor=north west
	] (router-services) {
	};
	
	% network interfaces
	\node[anchor=north] at (router-services.north east -| router-description.east) (router-interfaces) {
		\begin{tcbitemize}[interfaces raster]
			\tcbitem[remember as=router-eth0] eth0\hfill\faIcon{ethernet}
			\tcbitem[remember as=router-eth1] eth1\hfill\faIcon{ethernet}
		\end{tcbitemize}
		
	};
	
	% draw box around all parts defining a host
	\begin{pgfonlayer}{hostboxes}
		\node[
			host box,
			fit=(router-description.center) (router-services.south) (router-services.west) (router-interfaces.south)
		] (router-host-box) {};
	\end{pgfonlayer}
	% --------------- end router ---------------
	

	% --------------- start acid ---------------
	% host description
	\node[host description, below=3cm of router-host-box] (acid-description) {%
		\begin{tcolorbox}[
			adjusted title=Acid dosing PLC,
			host description box
		]
		IP: \hspace*{\fill} \ip{X.X.X.X}{subnet-home}
		\tcblower
		S7-1200
		\end{tcolorbox}
	};

	% services
	\node[
		below=0em of acid-description.south west,
		anchor=north west
	] (acid-services) {
		\begin{tcboxedraster}[services raster]{services box}
			\tcbox[enhanced, remember as=raspberrypi-pihole]{\service{}{Application}}
		\end{tcboxedraster}
	};

	
	% network interfaces
	\node[anchor=north] at (acid-services.north east -| acid-description.east) (acid-interfaces) {
		\begin{tcbitemize}[interfaces raster]
			\tcbitem[remember as=acid-eth0] \faIcon{ethernet}\hfill eth0
		\end{tcbitemize}
	};
	
	% draw box around all parts defining a host
	\begin{pgfonlayer}{hostboxes}
	\node[
		host box,
		fit=(acid-description.center) (acid-services.south) (acid-services.west) (acid-interfaces.south)
	] (acid-host-box) {};
	\end{pgfonlayer}
	% --------------- end acid ---------------
	

	% --------------- start raspberrypi ---------------
	% host description
	\node[host description, right=12cm of router-description] (raspberrypi-description) {%
		\begin{tcolorbox}[
			adjusted title= ScaleAQ Server,
			host description box
		]
		IP: \hspace*{\fill} \ip{X.X.X.X}{subnet-home}\\
		 SN: \hspace*{\fill} \ip{X.X.X.X}{subnet-home} \\
		Gateway: \hspace*{\fill} \ip{X.X.X.X}{subnet-home}\\
		VLAN: \hspace*{\fill} \ip{X}{subnet-home}
		\vspace{\baselineskip}
		
		\begin{tabular}{@{}ll@{}}
			\faIcon{compact-disc} & ESXi \\
		\end{tabular}
		
		\tcblower
		Supermicro 1U Server PC
		\end{tcolorbox}
	};
	
	% services
	\node[
		below=0em of raspberrypi-description.south east,
		anchor=north east
	] (raspberrypi-services) {
		\begin{tcboxedraster}[services raster]{services box}
			\tcbox[enhanced, remember as=raspberrypi-pihole]{\service{}{SiteControl}}
		\end{tcboxedraster}
	};

	
	% network interfaces
	\node[anchor=north] at (raspberrypi-services.north west -| raspberrypi-description.west) (raspberrypi-interfaces) {
		\begin{tcbitemize}[interfaces raster]
			\tcbitem[remember as=raspberrypi-eth0] \faIcon{ethernet}\hfill eth0
			\tcbitem[remember as=raspberrypi-eth1] \faIcon{ethernet}\hfill eth1
			\tcbitem[remember as=raspberrypi-eth2] \faIcon{ethernet}\hfill eth2
			\tcbitem[remember as=raspberrypi-eth2] \faIcon{ethernet}\hfill eth3
			\tcbitem[remember as=raspberrypi-eth2] \faIcon{ethernet}\hfill eth4
			\tcbitem[remember as=raspberrypi-eth2] \faIcon{ethernet}\hfill eth5
		\end{tcbitemize}
	};
	
	% draw box around all parts defining a host
	\begin{pgfonlayer}{hostboxes}
	\node[
		host box,
		fit=(raspberrypi-description.center) (raspberrypi-services.south) (raspberrypi-services.east) (raspberrypi-interfaces.south)
	] (raspberrypi-host-box) {};
	\end{pgfonlayer}
	% --------------- end raspberrypi ---------------
	

% --------------- begin liftUp ---------------
	% host description
	\node[host description, below=3cm of acid-host-box] (liftUp-description) {%
		\begin{tcolorbox}[
			adjusted title= LiftUp system,
			host description box
		]
		IP: \hspace*{\fill} \ip{X.X.X.X}{subnet-home}\\
		\vspace{\baselineskip}
		
		\tcblower
		PLC
		\end{tcolorbox}
	};
	
	% services
	\node[
		below=0em of liftUp-description.south west,
		anchor=north west
	] (liftUp-services) {
		\begin{tcboxedraster}[services raster]{services box}
			\tcbox[enhanced, remember as=liftUp-pihole]{\service{}{Stunner}}
			\tcbox[enhanced, remember as=liftUp-pihole]{\service{}{Counter camera}}
			\tcbox[enhanced, remember as=liftUp-pihole]{\service{}{Conveyor belt}}
			\tcbox[enhanced, remember as=liftUp-pihole]{\service{}{LiftUp}}
		\end{tcboxedraster}
	};

	
	% network interfaces
	\node[anchor=north] at (liftUp-services.north east -| liftUp-description.east) (liftUp-interfaces) {
		\begin{tcbitemize}[interfaces raster]
			\tcbitem[remember as=liftUp-eth0] \faIcon{ethernet}\hfill eth0
			\tcbitem[remember as=liftUp-eth1] \faIcon{ethernet}\hfill eth1
		\end{tcbitemize}
	};
	
	% draw box around all parts defining a host
	\begin{pgfonlayer}{hostboxes}
	\node[
		host box,
		fit=(liftUp-description.center) (liftUp-services.south) (liftUp-services.west) (liftUp-interfaces.south)
	] (liftUp-host-box) {};
	\end{pgfonlayer}
	% --------------- end LiftUp ---------------


	% --------------- begin barge equipment ---------------
	% host description
	\node[host description, below=3cm of liftUp-host-box] (barge-description) {%
		\begin{tcolorbox}[
			adjusted title= Barge control,
			host description box
		]
		IP: \hspace*{\fill} \ip{X.X.X.X}{subnet-home}\\
		\vspace{\baselineskip}
		
		\tcblower
		PLC
		\end{tcolorbox}
	};
	
	% services
	\node[
		below=0em of barge-description.south west,
		anchor=north west
	] (barge-services) {
		\begin{tcboxedraster}[services raster]{services box}
			\tcbox[enhanced, remember as=barge-pihole]{\service{}{Door control and monitoring}}
			\tcbox[enhanced, remember as=barge-pihole]{\service{}{Ventilation and heating}}
			\tcbox[enhanced, remember as=barge-pihole]{\service{}{Lighting system}}
			\tcbox[enhanced, remember as=barge-pihole]{\service{}{Cooling system}}
			\tcbox[enhanced, remember as=barge-pihole]{\service{}{Generator control and monitoring}}
			\tcbox[enhanced, remember as=barge-pihole]{\service{}{Switchboard / PMS}}
			\tcbox[enhanced, remember as=barge-pihole]{\service{}{Ballast system}}
			\tcbox[enhanced, remember as=barge-pihole]{\service{}{Bilge system}}
			\tcbox[enhanced, remember as=barge-pihole]{\service{}{Tank monitoring}}
			\tcbox[enhanced, remember as=barge-pihole]{\service{}{UPS monitoring}}
		\end{tcboxedraster}
	};

	
	% network interfaces
	\node[anchor=north] at (barge-services.north east -| barge-description.east) (barge-interfaces) {
		\begin{tcbitemize}[interfaces raster]
			\tcbitem[remember as=barge-eth0] \faIcon{ethernet}\hfill eth0
			\tcbitem[remember as=barge-eth1] \faIcon{ethernet}\hfill eth1
		\end{tcbitemize}
	};
	
	% draw box around all parts defining a host
	\begin{pgfonlayer}{hostboxes}
	\node[
		host box,
		fit=(barge-description.center) (barge-services.south) (barge-services.west) (barge-interfaces.south)
	] (barge-host-box) {};
	\end{pgfonlayer}
	% --------------- end barge equipment ---------------



		% --------------- begin fgs equipment ---------------
	% host description
	\node[host description, right=12cm of barge-description] (fgs-description) {%
		\begin{tcolorbox}[
			adjusted title= Fire and gas,
			host description box
		]
		IP: \hspace*{\fill} \ip{X.X.X.X}{subnet-home}\\
		\vspace{\baselineskip}
		
		\tcblower
		PLC
		\end{tcolorbox}
	};
	
	% services
	\node[
		below=0em of fgs-description.south east,
		anchor=north east
	] (fgs-services) {
		\begin{tcboxedraster}[services raster]{services box}
			\tcbox[enhanced, remember as=fgs-pihole]{\service{}{monitoring}}
			\tcbox[enhanced, remember as=fgs-pihole2]{\service{}{onboard alarm system}}
			\tcbox[enhanced, remember as=fgs-pihole3]{\service{}{alarm, system}}
			\tcbox[enhanced, remember as=fgs-pihole3]{\service{}{manual detectors}}
			\tcbox[enhanced, remember as=fgs-pihole3]{\service{}{signalling deviice}}
			\tcbox[enhanced, remember as=fgs-pihole3]{\service{}{emergency shutdown systems}}
		\end{tcboxedraster}
	};

	
	% network interfaces
	\node[anchor=north] at (fgs-services.north west -| fgs-description.west) (fgs-interfaces) {
		\begin{tcbitemize}[interfaces raster]
			\tcbitem[remember as=fgs-eth0] \faIcon{ethernet}\hfill eth0
			\tcbitem[remember as=fgs-eth1] \faIcon{ethernet}\hfill eth1
		\end{tcbitemize}
	};
	
	% draw box around all parts defining a host
	\begin{pgfonlayer}{hostboxes}
	\node[
		host box,
		fit=(fgs-description.center) (fgs-services.south) (fgs-services.east) (fgs-interfaces.south)
	] (fgs-host-box) {};
	\end{pgfonlayer}
	% --------------- end fgs ---------------


			% --------------- begin fgs equipment ---------------
	% host description
	\node[host description, right=12cm of fgs-description] (fwp-description) {%
		\begin{tcolorbox}[
			adjusted title= Fire and gas,
			host description box
		]
		IP: \hspace*{\fill} \ip{X.X.X.X}{subnet-home}\\
		\vspace{\baselineskip}
		
		\tcblower
		PLC
		\end{tcolorbox}
	};
	
	% services
	\node[
		below=0em of fwp-description.south east,
		anchor=north east
	] (fwp-services) {
		\begin{tcboxedraster}[services raster]{services box}
			\tcbox[enhanced, remember as=fwp-pihole]{\service{}{monitoring}}
			\tcbox[enhanced, remember as=fwp-pihole2]{\service{}{onboard alarm system}}
			\tcbox[enhanced, remember as=fwp-pihole3]{\service{}{alarm, system}}
			\tcbox[enhanced, remember as=fwp-pihole3]{\service{}{manual detectors}}
			\tcbox[enhanced, remember as=fwp-pihole3]{\service{}{signalling deviice}}
			\tcbox[enhanced, remember as=fwp-pihole3]{\service{}{emergency shutdown systems}}
		\end{tcboxedraster}
	};

	
	% network interfaces
	\node[anchor=north] at (fwp-services.north west -| fwp-description.west) (fwp-interfaces) {
		\begin{tcbitemize}[interfaces raster]
			\tcbitem[remember as=fwp-eth0] \faIcon{ethernet}\hfill eth0
			\tcbitem[remember as=fwp-eth1] \faIcon{ethernet}\hfill eth1
		\end{tcbitemize}
	};
	
	% draw box around all parts defining a host
	\begin{pgfonlayer}{hostboxes}
	\node[
		host box,
		fit=(fwp-description.center) (fwp-services.south) (fwp-services.east) (fwp-interfaces.south)
	] (fwp-host-box) {};
	\end{pgfonlayer}
	% --------------- end fgs ---------------

		

	\begin{scope}[ethernet]
		\draw[subnet-home] (raspberrypi-eth0.west) -- ++(-1,0) |- node[above,xshift = -6cm] {OPC UA} (router-eth0.east);
		\draw[subnet-home] (router-eth1.east) -- ++(1,0) |- (acid-eth0.east);
		\draw[subnet-home] (liftUp-eth0.east) -- ++(2,0) |- node[above,xshift = 1.5cm, yshift = -3cm] {MODBUS TCP/IP} (raspberrypi-eth1.west);
	\end{scope}


	\node[dashed, label=above:Automated dead fish handler, rounded corners, draw,inner sep=2cm, fit=(liftUp-services) (router-description) (acid-description)] (deadfish-box) {};'

	%\node[font=\large, above left=0cm and .5cm of deadfish-box] {Another text node.};

	% --------------- end network connections ---------------



\end{tikzpicture}


\end{document}